% RESEARCH_DRAFT_SHA256: 088e0287dabb9d4011eb7eec67b6ddaf0d42015a7f94c14db1f8e0d0319a97c4
% Layout and source-reference wrappers expanded for the web renderer.
\documentclass[11pt]{article}
\usepackage[margin=25mm]{geometry}
\usepackage{fontspec,unicode-math}
\setmainfont{Latin Modern Roman}
\setmathfont{Latin Modern Math}
\usepackage{amsmath,amssymb,mathtools,xurl,hyperref}
\hypersetup{colorlinks=true,urlcolor=blue,linkcolor=blue}
\setcounter{secnumdepth}{0}
\setlength{\parindent}{0pt}
\setlength{\parskip}{5pt}
\setlength{\emergencystretch}{3em}
\newcommand{\sref}[2]{\hyperlink{source:#1:#2}{[S#2]}}
\newcommand{\cataloguescope}[1]{\paragraph{Source record.}#1}
\newcommand{\floor}[1]{\left\lfloor#1\right\rfloor}
\begin{document}
\section*{Research-attempt provenance}
This standalone research attempt retains the complete original definition body below,
including its historical source and preparation notes. Those notes describe the
earlier catalogue preparation; the new attempt and its verification are identified
separately after that body. No original ID, catalogue status, or source statement is changed.

\section{A commuting dual-output contract for worksheet exercises}\label{statement}

\subsection{Definitions and mathematical statement}

\paragraph{Shared notation.}Standard number systems are $\mathbb N,\mathbb Z,\mathbb Q,\mathbb R,\mathbb C$. Specialized notation and index ranges follow the cited source and the definitions below.


PreTeXt represents document structure and mathematical content using semantic markup. A rendered output is a presentation produced from that source, such as web or print output. An exercise's mathematical prompt is distinct from its response interface and print layout. Writing workspace is physical blank area in print; an online interaction gathers a response through controls. A static realization presents an exercise without requiring its interactive execution.
\paragraph{Record scope and qualifications.}
The catalogue title is retained as metadata. Named certificates, contracts, bounds or frameworks in that title are not asserted as results of the original source extract.
The supplied sentence explicitly labels itself a reconstruction and is not a verified quotation. No exact semantic equivalence relation or required interaction class is supplied.

Catalogue screening limitation: PreTeXt software design, documentation, or publication feature request; no mathematical problem is specified by the supplied passage.

The source snapshot labels this record partially solved; this is catalogue metadata rather than independent verification of a mathematical result or the current completion of a project.
The source text below is rendered literally where necessary; damaged formulas and composite extracts are not silently repaired into a new mathematical claim.
The supplied catalogue marks this formulation as an unverified reconstruction; missing assumptions are not certified by the source link.
\paragraph{Statement recorded in the supplied source.}
Recovered design question (explicit reconstruction).** Can one semantic exercise source support (i) a static print realization that retains measured writing workspace and (ii) an ordinary online realization that retains the interactive response mechanism but suppresses paper-only blank space, while leaving a printable static realization of the interactive exercise available? \href{https://github.com/PreTeXtBook/community/wiki/Moderated-Problem-Session}{[S2]}
\paragraph{Source status.}
The supplied catalogue labels this source record partially solved; that metadata was checked against the supplied snapshot on 2026-10-11. The label is preserved for provenance and is not an independent verification of present open status or mathematical correctness.
\subsection{Short English statement}
Support print and interactive presentations from one exercise source while handling writing space appropriately.
\subsection{Sources}
\begin{itemize}
\item[S1] Ulam AI and the UnsolvedMath Contributors, UnsolvedMath: A Curated Collection of Open Mathematics Problems, catalogue entry. Catalogue curation and metadata: CC BY 4.0. The source extract is retained; definitions, English summaries and qualifications are editorial additions.
\url{https://huggingface.co/datasets/ulamai/UnsolvedMath}.
\item[S2] AIM Workshop Problem Lists, source linked by the catalogue.
\url{https://github.com/PreTeXtBook/community/wiki/Moderated-Problem-Session}.
\item[S3] Official PreTeXt Guide, Worksheets (context for worksheet and workspace terminology).
\url{https://pretextbook.org/doc/guide/html/overview-worksheet.html}.
\end{itemize}
\paragraph{Preparation and review.}
The original source and UnsolvedMath contributors retain attribution for the source extract. Codex (GPT-6; exact serving revision not exposed) assisted with definitions, English summaries, scope notes and formatting. No new human mathematical authorship or independently verified solution is asserted. Catalogue identity checked against the complete supplied registry and current repository. Distinct upstream records, including repeated questions and broad or damaged extracts, are retained for complete coverage of the source CS category. A repeated record is not a new independent conjecture. Complete mathematical-equivalence and current-literature review has not been performed. Preparation and final review follow contribution prompts C and D, with source gaps disclosed. Structural and rendering checks do not establish a new theorem or open status.
\label{end:statement}

\section{Literature and context check}
The community-session source was consulted directly. Its entry is a feature or design request, without a complete formal specification. The catalogue title is editorial and is not evidence of an existing theorem. The following finite model is an auxiliary investigation; no PreTeXt implementation or project completion is asserted. This record explicitly reconstructs a design question. The consulted worksheet guide already describes paper workspace that does not appear in HTML; that existing context is not claimed as a new result.
\section{Research attempt: Dual rendering and the irreversible loss of paper workspace}
Choose an auxiliary exercise representation $E=(q,A,v,h,s)$: prompt $q$, response set $A$, validator $v:A\to\{0,1\}$, paper workspace height $h\ge0$, and a static representation $s$ of the response mechanism. Define a paper rendering to retain $(q,s,h)$ and an interactive rendering to retain $(q,A,v)$, with a separate printable rendering available from the original $E$. Both retain the same exercise identifier and prompt. Under this declared model the two outputs are consistent on those shared observables by direct projection; this is a construction within the model, not a proof that a real renderer implements it.

The failed route is to generate the online representation first, discard paper-only data, and expect to recover the requested print output from that online representation alone. Let $E_1,E_2$ differ only in workspace heights $h_1\ne h_2$. Their online projections are identical but their paper projections differ. If a function $F$ of the online representation reconstructed paper output for every exercise, it would have to send that one shared input to two different outputs, a contradiction. Thus paper data must survive in the source or an attached side channel even when not displayed online.

An analogous obstruction applies to a screenshot fallback for response behavior. Let $A=\{0,1\}$ and two validators satisfy $v_1(a)=a$ and $v_2(a)=1-a$. A screenshot showing only the same prompt and empty answer box is identical for both. It cannot determine which answers are accepted. A printed static realization can faithfully present the task, but it does not preserve the interactive evaluation function unless that function is represented separately. Task content, writing space, and executable response behavior are therefore distinct observables.
\section{Verification and adversarial audit}
The argument is an exact analysis of the stated finite model. Its counterexample was checked by direct substitution; no PreTeXt build, browser test, or user study was run. The model's assumptions are explicit and are not asserted to describe every real document.
\paragraph{Audit corrections.} The existing worksheet behavior is recognized as prior functionality. Projection identities are not described as a implemented commuting renderer, and the static fallback is not claimed behaviorally equivalent to an interactive validator.
\paragraph{External computation.} None; symbolic arguments and boundary checks only.
\section{Scope and remaining gaps}
The calculation identifies a concrete failed route or a restricted correctness property. Actual schema integration, transformations, usability, and compatibility with the full original feature request remain unimplemented and unverified. Thus the original agenda remains unresolved, with no claimed mathematical novelty.
\paragraph{Integration review.} Codex (GPT-6; exact serving revision not exposed) reviewed the written arguments and expanded layout and source-reference wrappers. Repository metadata and the production web renderer were checked using Python 3.13 and JavaScript with MathJax 3.2.2. This model-assisted review is not independent human proof verification; no PDF compilation or formal proof checking is asserted.

\section{Completion Estimate}
\noindent\textbf{COMPLETION: 0\%}
The calculations below the preserved source statement investigate a concrete route,
but establish no substantiated new progress toward the entire original target.
Known special cases and failed sufficient conditions are not credited as a solution.
\section{References}
The original source references and catalogue attribution are preserved in the
Sources section above. The following URLs were actually consulted for this attempt;
retention of an additional original reference does not claim fresh consultation of it.
\begin{itemize}
\item \url{https://github.com/PreTeXtBook/community/wiki/Moderated-Problem-Session} (consulted 2026-10-11). 
\item \url{https://pretextbook.org/doc/guide/html/overview-worksheet.html} (consulted 2026-10-11). 
\end{itemize}
\section{Model and contributor disclosure}
OpenAI gpt-6-astra, requested reasoning configuration ultra, generated this research
attempt on 2026-10-11. The exact serving revision is not exposed. Codex GPT-6
performed orchestration, repository-format review, and a model-assisted mathematical
audit of the written arguments; its exact serving revision
is also not exposed. The mathematical route, boundary checks and final adversarial
self-audit followed contribution prompts A then B in
\texttt{docs/contribution\_prompts.md}, after reading \texttt{CONTRIBUTING.md}.
Python 3.14.2 assembled the standalone documents; any mathematical computation is
specified in Verification. No theorem prover was used. No human mathematical
authorship, substantial human proof correction, or independent human proof checking
is asserted. Model-assisted review, source retrieval, and structural checks do not
establish correctness or novelty. Full private reasoning transcripts are unavailable
as reproducibility artifacts; the complete arguments, inputs and outcomes actually
used are stated here or in the companion batch audit.
\section{Final status}
Unresolved as a new research contribution. This is a concrete route investigation,
not a claim of new complete resolution or a current-status certification. A known
answer reproduced for context does not become a new contribution. The catalogue's
source-status label is preserved; this attempt status does not assert that a
source-labelled solved question is open.
\end{document}
